%-------------------------
% Resume in Latex
% Author : Sourabh Bajaj
% Website: https://github.com/sb2nov/resume
% License : MIT
%------------------------

\documentclass[letterpaper,11pt]{article}

\usepackage[utf8]{inputenc}
\usepackage[T1]{fontenc}
\usepackage[brazil]{babel}
\usepackage{latexsym}
\usepackage[empty]{fullpage}
\usepackage{titlesec}
\usepackage{marvosym}
\usepackage[usenames,dvipsnames]{color}
\usepackage{verbatim}
\usepackage{enumitem}
\usepackage[pdftex]{hyperref}
\usepackage{fancyhdr}
\usepackage[sfdefault]{universalis}


\pagestyle{fancy}
\fancyhf{} % clear all header and footer fields
\fancyfoot{}
\renewcommand{\headrulewidth}{0pt}
\renewcommand{\footrulewidth}{0pt}

% Adjust margins
\addtolength{\oddsidemargin}{-0.375in}
\addtolength{\evensidemargin}{-0.375in}
\addtolength{\textwidth}{1in}
\addtolength{\topmargin}{-.5in}
\addtolength{\textheight}{1.0in}

\urlstyle{same}

\raggedbottom
\raggedright
\setlength{\tabcolsep}{0in}

% Sections formatting
\titleformat{\section}{
  \vspace{-4pt}\scshape\raggedright\large
}{}{0em}{}[\color{black}\titlerule \vspace{-5pt}]

%-------------------------
% Custom commands
\newcommand{\resumeItem}[2]{
  \item\small{
    \textbf{#1}{: #2 \vspace{-2pt}}
  }
}

\newcommand{\resumeSubheading}[4]{
  \vspace{-1pt}\item
    \begin{tabular*}{0.97\textwidth}{l@{\extracolsep{\fill}}r}
      \textbf{#1} & #2 \\
      \textit{\small#3} & \textit{\small #4} \\
    \end{tabular*}\vspace{-5pt}
}

\newcommand{\resumeSubItem}[2]{\resumeItem{#1}{#2}\vspace{-4pt}}

\renewcommand{\labelitemii}{$\circ$}

\newcommand{\resumeSubHeadingListStart}{\begin{itemize}[leftmargin=*]}
\newcommand{\resumeSubHeadingListEnd}{\end{itemize}}
\newcommand{\resumeItemListStart}{\begin{itemize}}
\newcommand{\resumeItemListEnd}{\end{itemize}\vspace{-5pt}}

%-------------------------------------------
%%%%%%  CV STARTS HERE  %%%%%%%%%%%%%%%%%%%%%%%%%%%%


\begin{document}

%----------HEADING-----------------
\begin{tabular*}{\textwidth}{l@{\extracolsep{\fill}}r}
  \textbf{\href{https://luisf.dev/}{\Large Luis Felipe Santos do Nascimento}} & Email : \href{mailto:luisfelipesdn12@gmail.com}{luisfelipesdn12@gmail.com}\\
  Desenvolvedor Full Stack Sênior $\cdot$ São Paulo, Brasil & Telefone : +55 11 95884-3078 \\
  \href{https://luisf.dev/}{https://luisf.dev} & \\
\end{tabular*}


%-----------PERFIL-----------------
\section{Perfil}
\small Desenvolvedor full stack sênior com base sólida em \textbf{TypeScript, Python e Go}. Experiência em APIs RESTful, automações, \textbf{web scraping}, integrações com modelos de \textbf{IA/LLMs}, interfaces internas e plataformas \textit{data-driven} em escala (milhões de dispositivos/usuários). Atua transformando problemas de negócio em soluções técnicas simples, escaláveis e funcionais, com mentalidade de melhoria contínua e colaboração.
\vspace{2pt}


%-----------FORMAÇÃO ACADÊMICA-----------------
\section{Formação Acadêmica}
  \resumeSubHeadingListStart
    \resumeSubheading
      {Universidade Presbiteriana Mackenzie}{São Paulo, Brasil}
      {Bacharel em Sistemas de Informação; Em andamento}{2023 -- 2027}
    \resumeSubheading
      {ETEC Basilides de Godoy}{São Paulo, Brasil}
      {Técnico em Desenvolvimento de Sistemas; Formado}{2021 -- 2022}
  \resumeSubHeadingListEnd


%-----------EXPERIÊNCIA-----------------
\section{Experiência}
  \resumeSubHeadingListStart

    \resumeSubheading
      {Painel TAP}{São Paulo, Brasil}
      {Desenvolvedor Full Stack Sênior}{mar. 2025 -- atualmente}
      \resumeItemListStart
        \resumeItem{Desenvolvimento Web Full-Stack}
          {Desenvolvimento full-stack de plataforma SaaS com Pyhton, Node, expondo APIs RESTful, interfaces internas e dashboards \textit{data-driven}. Automações, web scraping e integrações com modelos de IA/LLMs, com foco em arquitetura, performance e escalabilidade.}
      \resumeItemListEnd

    \resumeSubheading
      {Allcom Telecom}{São Paulo, Brasil}
      {Desenvolvedor Full Stack Pleno}{jun. 2024 -- mar. 2025}
      \resumeItemListStart
        \resumeItem{Plataforma IoT em escala (4M+ dispositivos, 8.000+ clientes)}
          {Atuação em produtos da maior operadora de IoT da América Latina, com parceiros estratégicos como \textbf{Claro, Vivo, Qualcomm e Nokia}. Desenvolvimento web full-stack com Next.js, Node, React, Laravel, Python, Docker e bancos relacionais. APIs RESTful, automações, web scraping e interfaces internas para segmentos de meios de pagamento (POS), smart cities, utilities, telemetria e frotas.}
      \resumeItemListEnd

    \resumeSubheading
      {Origin9}{São Paulo, Brasil}
      {Desenvolvedor Full Stack Pleno}{dez. 2023 -- jun. 2024}
      \resumeItemListStart
        \resumeItem{Desenvolvimento Web Full-Stack}
          {Desenvolvimento web front e back-end com Node, Python, React e Docker. APIs RESTful, automações, web scraping e bancos relacionais sob metodologias ágeis.}
      \resumeItemListEnd

    \resumeSubheading
      {YOUDevelop}{São Paulo, Brasil}
      {Desenvolvedor Full Stack Pleno}{out. 2021 -- dez. 2023}
      \resumeItemListStart
        \resumeItem{Desenvolvimento Web Full-Stack}
          {Desenvolvimento web front e back-end com Next, React, Firebase e Supabase. Construção de APIs RESTful, integrações de dados e entrega contínua em metodologias ágeis.}
      \resumeItemListEnd

    \resumeSubheading
      {Morph}{São Paulo, Brasil}
      {Full Stack Web \& Mobile Engineer}{dez. 2021 -- jun. 2022}
      \resumeItemListStart
        \resumeItem{Desenvolvimento Mobile}
          {Desenvolvimento ágil em Flutter para aplicações Web e Mobile.}
      \resumeItemListEnd

    \resumeSubheading
      {Zyndicate}{São Paulo, Brasil}
      {Full Stack Engineer}{dez. 2020 -- dez. 2021}
      \resumeItemListStart
        \resumeItem{Desenvolvimento Web}
          {Desenvolvimento Web utilizando GitHub Flow, ReactJS, NextJS, GraphQL, UI/UX, CI/CD, Metodologias Ágeis e conversação em inglês.}
      \resumeItemListEnd

    \resumeSubheading
      {Nooio}{São Paulo, Brasil}
      {Backend Developer}{dez. 2020 -- mai. 2021}
      \resumeItemListStart
        \resumeItem{Automações, Scraping e APIs}
          {Desenvolvimento de APIs RESTful com Node/Express e scripts em Python para automações e web scraping. Desenvolvimento Guiado por Testes (TDD) e deploy em AWS.}
      \resumeItemListEnd

  \resumeSubHeadingListEnd


%-----------CURSOS-----------------
\section{Cursos}
  \resumeSubHeadingListStart
    \resumeSubItem{Linguagem de Programação Go -- Aprenda Go (Ellen Körbes)}
      {Estruturas de dados, interfaces e polimorfismo, concorrência e paralelismo, tratamento de erros, testes e benchmarks. 2020 -- 22h/aula.}
    \resumeSubItem{Curso de Python 3 -- Curso em Vídeo (Gustavo Guanabara)}
      {Fundamentos sólidos da linguagem, de variáveis a tratamento de erros. 2020 -- 120h/aula.}
    \resumeSubItem{Curso Backend Developer -- Digital Innovation One}
      {Arquitetura de Sistemas, Java, Spring Boot, testes, SQL, Node.js com Express, melhores práticas com APIs e LGPD. 2020 -- 100h/aula.}
  \resumeSubHeadingListEnd


%-----------CERTIFICAÇÕES-----------------
\section{Certificações}
  \resumeSubHeadingListStart
    \resumeSubItem{Docker: criando e gerenciando containers}{Fundamentos de containers, imagens, volumes e orquestração com Docker.}
    \resumeSubItem{Bootcamp Backend Developer Carrefour}{Programa intensivo de desenvolvimento backend com foco em Java, Spring Boot e melhores práticas de mercado.}
    \resumeSubItem{Desenvolvimento back-end com Node.js}{Construção de APIs e serviços com Node.js e Express.}
    \resumeSubItem{Desenvolvimento avançado com JavaScript ES6}{Recursos modernos do JavaScript (ES6+), padrões e melhores práticas.}
    \resumeSubItem{Expert Lesson: Introdução à LGPD}{Princípios da Lei Geral de Proteção de Dados e impactos no desenvolvimento de software.}
  \resumeSubHeadingListEnd


%-----------RECONHECIMENTOS-----------------
\section{Reconhecimentos}
  \resumeSubHeadingListStart
    \resumeSubItem{Finalista da OBI (Olimpíada Brasileira de Informática)}{Classificado para a fase final da competição nacional de algoritmos e programação.}
    \resumeSubItem{Melhor aluno do curso de Desenvolvimento de Sistemas}{Reconhecimento de melhor aluno da turma técnica em Desenvolvimento de Sistemas -- ETEC Basilides de Godoy.}
  \resumeSubHeadingListEnd

%
%--------HABILIDADES------------
\section{Habilidades}
 \resumeSubHeadingListStart
   \item{
     \textbf{Linguagens}{: Python, Go, Node.js, TypeScript, JavaScript (ES6+), Java, Dart, Arduino}
   }
   \item{
     \textbf{Dados, IA \& Scraping}{: Grafana, Pandas, BeautifulSoup, Scrapy, LLMs / IA aplicada, Data Visualization, bancos relacionais (SQL), Firebase, Supabase}
   }
   \item{
     \textbf{Web \& APIs}{: APIs RESTful, GraphQL, Next.js, React, Express, Laravel, Flutter, UI/UX, Acessibilidade digital}
   }
   \item{
     \textbf{Infra \& Cloud}{: AWS, Docker, CI/CD, Git/GitHub, Arquitetura de Sistemas}
   }
   \item{
     \textbf{Práticas}{: TDD, Automação, Metodologias Ágeis (SCRUM), LGPD}
   }
   \item{
     \textbf{Idiomas}{: Português (nativo), Inglês (avançado -- Full Professional)}
   }
 \resumeSubHeadingListEnd


%-------------------------------------------
\end{document}
